\documentclass[10pt]{article}
\usepackage{amsmath} % aligned
\usepackage{geometry}
\geometry{margin=1in}
\newcommand{\bfL}{{\bf L}}
\newcommand{\bfA}{{\bf A}}
\newcommand{\bfb}{{\bf b}}
\newcommand{\bfB}{{\bf B}}
\newcommand{\bfx}{{\bf x}}
\newcommand{\bfv}{{\bf v}}
\newcommand{\bfna}{{\mbox{\boldmath $\bf\nabla$}}}

\begin{document}

\title{\bf Magnetic field, its derivatives, and the null point}
\date{}
\maketitle

\section{Biot-Savart field of a circular loop}

Biot-Savart field equals
\begin{equation}
\bfB(\bfx)=\oint\bfb(\bfx,t)\,dt,\quad
\bfb(\bfx,t) = \frac{\dot\bfL\times(\bfx-\bfL)}{D^3},\quad D\equiv|\bfx-\bfL|.
\end{equation}
Here the current loop is defined parametrically by $\bfL(t)=(a\cos(t),a\sin(t),0)$.

Component-wise, we have
\begin{equation}
\begin{aligned}
b_x &= \frac{\dot L_y(z-L_z) - \dot L_z(x-L_x)}{D^3} = \frac{\dot L_yz}{D^3},\\
b_y &= \frac{\dot L_z(x-L_x) - \dot L_x(z-L_z)}{D^3} = -\frac{\dot L_xz}{D^3},\\
b_z &= \frac{\dot L_x(y-L_y) - \dot L_y(x-L_x)}{D^3}
= \frac{L_x\dot L_y-L_y\dot L_x + \dot L_xy - \dot L_y x}{D^3}
= \frac{a^2 + \dot L_xy - \dot L_y x}{D^3}.
\end{aligned}
\end{equation}
Here the condition $L_z\equiv0$ is applied for a flat $(x,y)$ loop.  In terms of the parameter $t$,
\begin{equation}
\begin{aligned}
b_x &= \frac{az\cos(t)}{D^3},\\
b_y &= \frac{az\sin(t)}{D^3},\\
b_z &= a\frac{a - x\cos(t) - y\sin(t)}{D^3},\\
D   &= \sqrt{(x - a\cos(t))^2 + (y - a\sin(t))^2 + z^2}
\end{aligned}
\end{equation}
Note that, for $y=0$, $D$ is an even function of $t$, and $B_y=\int_{-\pi}^\pi b_y\,dt=0$.

Derivatives of the infinitesimal field:
\begin{equation}
\begin{aligned}
b_{x,x} &= -3az\frac{\cos(t)(x-a\cos(t))}{D^5},\\
b_{y,y} &= -3az\frac{\sin(t)(y-a\sin(t))}{D^5},\\
b_{z,z} &=  3az\frac{x\cos(t) + y\sin(t) - a}{D^5},\\
b_{x,z} &=  \frac{a\cos(t)}{D^3}\left(1-3\frac{z^2}{D^2}\right),\\
b_{z,x} &= -\frac{a\cos(t)}{D^3} -3a\frac{(x-a\cos(t))(a-x\cos(t)-y\sin(t))}{D^5}.
\end{aligned}
\end{equation}
It is immediately verified that
$\bfna\cdot\bfb=b_{x,x}+b_{y,y}+b_{z,z}=0$, and therefore the total
field $\bfB$ is divergence free, as it should be.  Not so for the
infinitesimal $b_{x,z}$ and $b_{z,x}$.  The current-free condition
$B_{x,z}=B_{z,x}$ is valid only for the total field after $\oint dt$
over $[0,2\pi]$.  Proving this in general by integrating the
derivatives above by $t$ is very difficult, as elliptic integrals are
involved.  {\sl Mathematica} consistenly evaluates $B_{x,z}-B_{x,z}$ at
various random values of $(a,x,y,z)$ as exact zero, which cannot be a
coincidence.

\subsection{{\sl Mathematica} checks}

Wolfram Cloud was used to verify above derivative formulas.  Some of
{\tt TeXForm} expressions:
\begin{equation}
\begin{aligned}
b_{x,z} &= \frac{a \cos (t) \left(a^2-2 a x \cos (t)-2 a y \sin (t)+x^2+y^2-2
   z^2\right)}{\left((x-a \cos (t))^2+(y-a \sin (t))^2+z^2\right)^{5/2}},\\
b_{z,x} &= \frac{a\left(-\cos(t)\left(-2a^2+ay\sin(t)-2x^2+y^2+z^2\right)-3x(a-y\sin(t))-ax\cos^2(t)\right)}{\left(a^2-2ax\cos(t)-2ay\sin(t)+x^2+y^2+z^2\right)^{5/2}}.
\end{aligned}
\end{equation}

\section{The null ("$X$") point}

The field in the plane $y=0$ is expressed through the vector potential
$A$ as $(B_x, B_z) = (-\partial_zA_y, \partial_xA_y)$, meaning that
the countour lines of $A_y(x, z)$ are the magnetic lines, and the $X$
point is a saddle point of $A_y$.  The $X$ point condition, however,
can be written in terms of $(B_x,B_z)$ only.  The field is zero at the
target point $(X,Z)$ and its Taylor expansion to first order is
\begin{equation}
  \begin{aligned}
B_x &= B_{x,x}(x-X)+B_{x,z}(z-Z),\\
B_z &= B_{z,x}(x-X)+B_{z,z}(z-Z).
  \end{aligned}
  \label{linear}
\end{equation}
The commas denote differentiation.  Unlike the field $\bfB$, its
derivatives are not zero at the $X$ point.

Magnetic lines near the $X$ point are either hyperbolas or
intersecting separatrices passing through the point.  The separatrices
are the directions where the field is collinear with the displacement,
$B_z/B_x = \pm(z-Z)/(x-X)$.  The desired direction is
$T=(z-Z)/(x-X)=\tan(36^\circ)$, where $T$ is the tangent of the
required angle between the field and the board plane ($54^\circ$ with
respect to the normal).  The collinearity condition yields
\begin{equation}
\frac{B_{z,x} + B_{z,z}T}{B_{x,x}+B_{x,z}T} = \pm T.
\end{equation}
For curl-free field, we have $B_{z,x} = B_{x,z}=0$, so the angle
condition involves only 3 derivatives.  For the plus sign above, the
slope of the separatrix is determined from the quadratic equation
\begin{equation}
B_{x,z}T^2 + (B_{x,x} - B_{z,z})T - B_{z,x} = 0\quad\Longrightarrow\quad
T=\frac{C+\sqrt{C^2+4B_{x,z}B_{z,x}}}{2B_{x,z}},\quad C\equiv B_{z,z}-B_{x,x}.
\end{equation}
This expression agrees with the fact [known to AI] that the separatrix
lines are the eigenvectors of the feld's Jacobian matrix:
\begin{equation}
J=\left(\begin{array}{ll} B_{x,x} & B_{x,z}\\ B_{z,x} & B_{z,z}\end{array}\right),\quad J\bfv=\lambda\bfv,\quad
\bfv_{1,2}=\left(C\pm\sqrt{C^2+4B_{x,z}B_{z,x}}, 2B_{x,z}\right).
\end{equation}
Note that the two eignevectors are orthogonal: $\bfv_1\cdot\bfv_2=0$
due to the curl-free field, $B_{x,z}=B_{z,x}$, as expected.  The
expression for the eigenvectors in {\tt mark\_x()} is used for a visual
check of the $X$ point separatrix relative to magnetic field lines.

\section{Gradient}

Given that the direction $\bfv$ of the separatrix near the null point
is collinear with the direction of $\bfB$, the gradient of
$|\bfB|=(\bfv\cdot\bfB)/|\bfv|$ along $\bfv$ can be written as
\begin{equation}
|\bfna B| = |\bfv|^{-2}(\bfv\cdot\bfna)(\bfv\cdot\bfB).
\end{equation}
Using the linear expansion of the components of $\bfB$ in
(\ref{linear}), and the fact that $B_{x,z}=B_{z,x}$, the gradient is
\begin{equation}
|\bfna B|=\frac{v_x^2B_{x,x} + 2v_xv_zB_{x,z}+v_z^2B_{z,z}}{v_x^2+v_z^2}.
\end{equation}

\section{Optimization of magnetic field}

We are concerned with a 2D magnetic geometry in the $(x,z)$ plane.



\end{document}
